% Created 2026-03-18 Wed 03:02
% Intended LaTeX compiler: lualatex
\documentclass[11pt]{article}
\usepackage{amsmath}
\usepackage{fontspec}
\usepackage{graphicx}
\usepackage{longtable}
\usepackage{wrapfig}
\usepackage{rotating}
\usepackage[normalem]{ulem}
\usepackage{capt-of}
\usepackage{hyperref}
\author{Prologos Project}
\date{2026-03-18}
\title{Gödel Completeness: Decidability as an Architectural Property}
\hypersetup{
 pdfauthor={Prologos Project},
 pdftitle={Gödel Completeness: Decidability as an Architectural Property},
 pdfkeywords={},
 pdfsubject={},
 pdfcreator={Emacs 30.2 (Org mode 9.7.39)}, 
 pdflang={English}}
\begin{document}

\maketitle
\tableofcontents

\section{Purpose}
\label{sec:orgc38d131}

This document codifies the principle that all computation hosted on
Prologos's propagator networks should be provably decidable ---
terminating for all inputs, with the termination argument stated
explicitly for each layer and propagator. We call this property
\textbf{Gödel Completeness}: the dual of Turing Completeness, where every
computation provably halts rather than potentially diverging.

It is far more difficult to design a decidable computation model than
to accidentally exercise the Halting Problem. Most systems default to
Turing Completeness and add fuel bounds as a safety net. Prologos
inverts this: the propagator network is a \textbf{decidable kernel} where
termination is structural, and general recursion lives outside it.

This principle emerged from the Track 6 propagator migration, where
the design of stratified propagator networks (S(-1), S0, L1, L2) raised
the question: what guarantees does layered quiescence provide? The
answer led to a hierarchy of termination guarantees and a discipline
for applying them.

For how this principle interacts with propagator-first infrastructure
and correct-by-construction design, see \href{DESIGN\_PRINCIPLES.org}{DESIGN\textsubscript{PRINCIPLES.org}}.
\section{The Decidability Principle}
\label{sec:org52bf8c2}

\begin{quote}
Every computation hosted on the propagator network must come with a
termination guarantee --- structural proof preferred, fuel bound as
pragmatic fallback. The network is a decidable computation model;
general recursion and unbounded computation happen outside it. When a
new propagator or lattice is added, its termination argument must be
stated explicitly.
\end{quote}

This is a design discipline, not a blanket theorem. It is a commitment
to \emph{thinking about termination} at every level, and choosing the
strongest guarantee available. The fuel fallback ensures pragmatism;
the discipline ensures we don't default to fuel when a structural
argument exists.
\subsection{What "Gödel Complete" means}
\label{sec:org279927d}

A system is Gödel Complete when every well-formed computation in the
system provably terminates. This is the opposite of Turing
Completeness, which guarantees expressiveness at the cost of
undecidable halting.

\begin{center}
\begin{tabular}{lll}
Property & Turing Complete & Gödel Complete\\
\hline
Expressiveness & All computable fns & Total computable fns\\
Halting & Undecidable & Decidable (structural)\\
Paradigm & Partial functions & Total functions\\
Examples & Haskell, Racket & Agda, Idris, Walnut\\
Prologos analogue & User-level \texttt{defn} & Propagator network\\
\end{tabular}
\end{center}

The propagator network is the Gödel Complete kernel. User-level code
(\texttt{defn} with general recursion, IO effects) is Turing Complete. The
boundary is clear: the network reasons \emph{about} programs; programs
\emph{execute} outside the network.
\section{The Hierarchy of Termination Guarantees}
\label{sec:org1adf136}

Not all termination arguments are equal. We recognize five levels,
ordered from strongest (most informative) to weakest (least
informative). The discipline is: use the strongest level available.
\subsection{Level 1: Structural (Tarski Fixpoint)}
\label{sec:orge82c92e}

Finite lattice + monotone propagators → fixpoint in at most \emph{h} steps,
where \emph{h} is the lattice height. No fuel needed. The lattice height IS
the bound.

\textbf{When it applies}: Any propagator operating on a finite domain with a
monotone merge function.

\textbf{Examples in Prologos}:
\begin{itemize}
\item S0 type propagation: each meta cell transitions at most once (⊥ →
solved). Finite metas = finite writes.
\item Effect propagation: effect lattice has finite height (declared
effects are finite).
\item QTT multiplicity: 3-element lattice (\texttt{m0 < m1 < mw}).
\item S(-1) retraction: believed-assumptions set can only shrink. Finite
assumptions → finite retractions.
\item L1 readiness: each readiness propagator fires at most once per
dependency becoming non-⊥.
\end{itemize}

\textbf{Proof obligation}: Show that the lattice has finite height and the
propagator's transfer function is monotone.
\subsection{Level 2: Well-Founded Measure}
\label{sec:orgd274b95}

Cross-stratum feedback or recursive resolution decreases a well-
founded ordering on each step. Termination follows from well-
foundedness (every strictly decreasing chain in a well-founded order
is finite).

\textbf{When it applies}: When one propagator layer's output triggers another
layer, and the triggering chain has a natural decreasing measure.

\textbf{Examples in Prologos}:
\begin{itemize}
\item L2→S0 feedback (trait resolution): resolving a trait constraint may
create new metas and new constraints, but at strictly smaller type
depth. Type depth is a natural number (well-founded). Each cycle
decreases depth → finite cycles.
\item Inductive structural recursion: computation over Nat, List, Tree
decreases the structural size on each recursive step.
\item Narrowing (functional-logic sense): each narrowing step reduces the
set of possible values --- the search space shrinks on a well-
founded measure.
\end{itemize}

\textbf{Proof obligation}: Identify the measure, show it decreases on every
step, show the ordering is well-founded.

\textbf{Enforcement}: For trait resolution, the decreasing type-depth
condition can be checked at instance registration time (correct by
construction). Instances that violate the condition are rejected with
a diagnostic --- no fuel needed. This is analogous to structural
recursion checking in Agda/Idris.
\subsection{Level 3: Widening-Guaranteed Convergence}
\label{sec:orge06f5b6}

The concrete domain is infinite, but a widening operator forces
convergence to an abstract fixpoint in bounded steps. Precision is
traded for termination. Narrowing (lattice sense) then recovers
precision without threatening termination (bounded by the number of
widening steps).

\textbf{When it applies}: Abstract interpretation over infinite domains
(integers, sizes, resource bounds, numeric ranges).

\textbf{Examples in Prologos}:
\begin{itemize}
\item Integer range analysis: concrete domain is ℤ (infinite height).
Widen to sign domain \{neg, zero, pos, top\} (height 2). Fixpoint
in at most 2 steps per cell.
\item Resource bound analysis: concrete usage counts are unbounded. Widen
to multiplicity lattice \{m0, m1, mw\} (height 2).
\end{itemize}

\textbf{Proof obligation}: Show that the widening operator produces a lattice
of finite height, and that narrowing is bounded by widening steps.

\textbf{Relationship to Level 1}: After widening, the abstract domain IS a
finite lattice --- so Level 3 reduces to Level 1 on the abstract
domain. The novelty is that the \emph{concrete} domain is infinite; the
widening operator is what makes it tractable.
\subsection{Level 4: Coinductive Productivity}
\label{sec:org46cfe20}

The computation is infinite (a stream, a session, an ongoing process),
but each \emph{observation} terminates. Productivity replaces termination
--- the guarantee is "you can always get the next element," not "the
computation eventually stops."

This is the greatest fixpoint (gfp) dual of inductive termination
(lfp). Where induction says "the base case is reached in finite
steps," coinduction says "each productive step completes in finite
time, yielding the next layer of the infinite structure."

\textbf{When it applies}: Stream processing, session types (ongoing
communication), reactive systems, coinductive types.

\textbf{Examples in Prologos}:
\begin{itemize}
\item Session types: a session type like \texttt{!Int.?Bool.!Int.?Bool...} is
coinductive (potentially infinite). But each message
send/receive is a finite observation. The propagator network
computes one "layer" of the session per quiescence cycle.
\item Stream processing: \texttt{Stream Int} is coinductive. Computing \texttt{head}
terminates; computing \texttt{tail} yields another stream (productive
observation, not full evaluation).
\item Effect system bridges: the IO bridge processes one effect
description per cycle. The stream of effects may be unbounded,
but each effect's processing terminates.
\end{itemize}

\textbf{Proof obligation}: Show that each observation (one layer of the
coinductive structure) can be computed in finitely many propagator
steps. This often reduces to Level 1 or 2 for the individual layer.

\textbf{Relationship to the propagator network}: Coinductive computations
on the network use \emph{productive quiescence} --- the network reaches
quiescence for one layer (terminating), yields an observation, then
is perturbed by the next input (the next layer begins). The per-
layer quiescence is Level 1; the overall computation is productive
(Level 4).
\subsection{Level 5: Fuel Bound}
\label{sec:org3df24c0}

Decidable by fiat. The computation stops because we force it to after
\emph{n} steps. Informative about safety (the system cannot diverge), not
about structure (we don't know WHY it terminates within the bound).

\textbf{When it applies}: As defense-in-depth behind Levels 1-4. As the
primary guarantee only when no structural argument is available (e.g.,
user-opted-into undecidable instance resolution).

\textbf{Examples in Prologos}:
\begin{itemize}
\item Stratified resolution fuel (100 iterations): currently the primary
guarantee. Track 7 will add Level 2 (well-founded measure) as the
structural guarantee, with fuel retained as defense-in-depth.
\item Reduction fuel: user-level computation (\texttt{defn} with general
recursion) runs with fuel. This is correct --- user code IS
potentially non-terminating; fuel is the appropriate mechanism.
\end{itemize}

\textbf{Role in the architecture}: Fuel should be the \emph{safety net}, not the
\emph{design}. When a computation relies on fuel as its primary termination
guarantee, ask: "can we identify a structural argument (Levels 1-4)
that makes the fuel bound unnecessary for correctness, retaining it
only for defense-in-depth?"
\section{Infinite Domains on a Decidable Network}
\label{sec:org3d10aeb}

A common misconception: propagator networks are limited to finite
domains. Through the combination of inductive/coinductive reasoning,
abstract interpretation, and widening/narrowing, the Prologos
propagator network can handle infinite domains while maintaining
decidability.
\subsection{Inductive Domains (lfp)}
\label{sec:orgc0cdae2}

Nat, List, Tree --- finite structures built from constructors. Propagation
over these terminates because each inductive value has finite depth.
The structural recursion principle gives a well-founded measure
(Level 2). The propagator network handles inductive types naturally:
each constructor application is a cell write; pattern matching is a
propagator that reads the cell and dispatches.
\subsection{Coinductive Domains (gfp)}
\label{sec:org58bcc09}

Streams, sessions, processes --- potentially infinite structures observed
one layer at a time. The propagator network computes one productive
observation per quiescence cycle (Level 4). The per-observation
computation is Level 1 (finite lattice within the cycle). The
combination gives decidability for each step and productivity for the
overall computation.
\subsection{Abstract Domains (Galois Connection)}
\label{sec:org8730e0f}

For domains like ℤ or ℝ where the concrete lattice has infinite
height, abstract interpretation provides a finite abstraction. The
Galois connection (α, γ) between concrete and abstract ensures
soundness:

\begin{itemize}
\item α (abstraction): concrete → abstract (e.g., 42 → pos)
\item γ (concretization): abstract → concrete set (e.g., pos → \{1, 2, \ldots{}\})
\end{itemize}

Propagators operate on the abstract domain (Level 1 termination).
Widening forces convergence when even the abstract domain has
unbounded ascending chains. Narrowing (lattice sense) recovers
precision post-widening, bounded by widening steps.
\subsection{The IO Bridge: Effects as Pure Values}
\label{sec:org7129047}

The IO bridge propagator network captures effects as first-class
values --- descriptions of intent, not executed side effects. This is
structurally a graded monad:

\begin{itemize}
\item The effect annotation (the grade) is a lattice value propagated
through the network
\item \texttt{bind} sequences effect descriptions; the grade combines via the
effect quantale's tensor product
\item The propagator network is the "evaluator" operating on monadic values
\end{itemize}

Because effect descriptions are \emph{values} in the network (not side
effects), they are subject to the same termination guarantees as any
other cell value. The effect lattice has finite height (declared
effects are finite), so effect propagation terminates structurally
(Level 1). Only the \emph{execution} of effects (actual IO) happens
outside the network, in the imperative world.

This means we can reason about effectful computations \emph{within} the
decidable propagator framework. Effect composition, ordering, and
constraint checking all terminate structurally. The monadic bridge
brings effects into the decidable kernel without compromising
decidability.
\section{Application to the Stratified Propagator Network}
\label{sec:org8713e36}

The Track 7 stratified propagator network architecture assigns a
termination guarantee level to each layer:

\begin{center}
\begin{tabular}{llrl}
Layer & Role & Level & Argument\\
\hline
S(-1) & Retraction & 1 & Finite assumptions, monotone ↓\\
S0 & Type propagation & 1 & Finite metas, each solved once\\
L1 & Readiness detection & 1 & Finite constraints, fire once\\
L2 & Resolution commitment & 2 & Type depth decreases per cycle\\
Cross & L2→S0 feedback & 2 & Finite metas bounds cycle count\\
Full & Layered quiescence & 1+2 & Composition of per-layer guarantees\\
\end{tabular}
\end{center}

The full layered quiescence terminates because:
\begin{enumerate}
\item Each individual layer terminates (Level 1 or 2)
\item Cross-layer feedback (L2→S0) decreases a well-founded measure
(number of unsolved metas × maximum type depth)
\item The product of two well-founded orderings is well-founded
\end{enumerate}

Fuel is retained as defense-in-depth for the full cycle, catching
bugs in the termination argument. But the structural argument (not
the fuel) is what guarantees termination under correct operation.
\subsection{Trait Resolution Decidability}
\label{sec:org52b0724}

The potential source of non-termination in L2 is trait resolution
creating new metas. The decidability condition:

\begin{quote}
An instance declaration \texttt{impl Foo [T args...]} where \texttt{[Constraint deps...]}
is admissible iff every constraint's type depth is strictly less than
the head's type depth.
\end{quote}

This is checked at instance registration time (correct by
construction). Admissible instances guarantee that resolution chains
decrease in type depth on every step --- a well-founded measure
(Level 2).

For advanced users who need non-decreasing instances (analogous to
Haskell's \texttt{UndecidableInstances}), a pragma opts into fuel-bounded
resolution (Level 5). The user makes an explicit choice to leave
the decidable fragment. The pragma should be rare; the default is
decidability.
\section{The Decidability Boundary}
\label{sec:orge9b4ff4}

The propagator network is the decidable kernel. User-level code is
the Turing Complete shell. The boundary is architecturally clear:

\begin{center}
\begin{tabular}{lll}
Domain & Decidable? & Mechanism\\
\hline
Type propagation & Yes & Level 1 (Tarski)\\
Effect propagation & Yes & Level 1 (Tarski)\\
QTT multiplicity & Yes & Level 1 (Tarski)\\
Constraint readiness & Yes & Level 1 (Tarski)\\
Retraction & Yes & Level 1 (finite ↓)\\
Trait resolution & Yes* & Level 2 (depth ↓)\\
Abstract interpretation & Yes & Level 3 (widening)\\
Session type checking & Yes & Level 4 (productive)\\
Stream processing & Yes & Level 4 (productive)\\
User-level \texttt{defn} & No & Fuel (Level 5)\\
General IO & No & Outside network\\
\end{tabular}
\end{center}

(*) Decidable with the type-depth admissibility condition. Opt-out
pragma falls back to fuel.

The principle does not attempt to make user code decidable (that
would cripple expressiveness). It makes the \emph{infrastructure} decidable
--- the type checker, the constraint solver, the effect system, the
session type checker. These are the components that users depend on
to be reliable, predictable, and non-surprising. A type checker that
diverges is a broken tool; a user function that diverges is the
user's choice.
\section{Design Discipline}
\label{sec:org774b9a7}

When adding a new propagator, lattice, or computation to the network:

\begin{enumerate}
\item \textbf{State the termination argument explicitly.} "This propagator
terminates because\ldots{}" One sentence, in the code comment or
design doc.

\item \textbf{Identify the guarantee level.} Which of Levels 1-5 applies?
Document it.

\item \textbf{If Level 5 (fuel), ask: can we do better?} Is there a structural
argument (Levels 1-4) that makes fuel unnecessary for correctness?
If yes, implement the structural guarantee and retain fuel as
defense-in-depth. If no, document why --- what property of the
domain prevents a structural argument?

\item \textbf{For cross-layer interactions, identify the well-founded measure.}
If layer A's output triggers layer B, which triggers layer A
again, what decreases on each cycle? If nothing decreases,
the cycle may not terminate --- this requires either a structural
fix or explicit fuel.

\item \textbf{For infinite domains, identify the finitization mechanism.}
Widening? Abstraction? Coinductive productivity? Document
which mechanism makes the infinite domain tractable.
\end{enumerate}

This is not bureaucratic overhead --- it is the \emph{minimum viable
reasoning} for ensuring the network remains decidable as it grows.
A single propagator without a termination argument can make the
entire network potentially non-terminating.
\section{Relationship to Other Principles}
\label{sec:org30154ba}

\begin{itemize}
\item \textbf{Correct by Construction} (\href{DESIGN\_PRINCIPLES.org}{DESIGN\textsubscript{PRINCIPLES.org}}): Decidability IS
a correct-by-construction property. A decidable network cannot
diverge --- not because of fuel, but because of structural properties.
Gödel Completeness extends correct-by-construction from "wrong
states are unrepresentable" to "non-termination is unrepresentable."

\item \textbf{Propagator-First Infrastructure} (\href{DESIGN\_PRINCIPLES.org}{DESIGN\textsubscript{PRINCIPLES.org}}): The
propagator network's decidability gives it a trustworthiness that
imperative infrastructure lacks. A propagator-first system that is
also Gödel Complete provides deterministic, terminating responses
to every query. This is essential for tooling (LSP, IDE) where
divergence means a frozen editor.

\item \textbf{Stratified Propagator Networks} (\href{DESIGN\_PRINCIPLES.org}{DESIGN\textsubscript{PRINCIPLES.org}}): Stratification
is the mechanism that makes non-monotone operations (retraction,
aggregation) decidable. By containing non-monotonicity within a
stratum that runs to fixpoint before higher strata observe it,
stratification preserves the decidability of the overall network.
Retraction is aggregation's dual; both are addressed by stratification.

\item \textbf{Data-Oriented Design}: Termination arguments are \emph{data} about
propagators, not runtime checks. The lattice height, the well-
founded measure, the widening operator --- these are structural
properties declared at definition time, not guards checked at
execution time.
\end{itemize}
\section{Historical Context}
\label{sec:orgb51ba23}

The decidability principle emerged from the Track 6 propagator
migration (2026-03-17). The design of stratified propagator networks
raised the question of termination guarantees for layered quiescence.
This led to:

\begin{enumerate}
\item The hierarchy of termination guarantees (Levels 1-5)
\item The recognition that retraction is aggregation's dual, both
addressed by stratification
\item The connection between the IO bridge's monadic structure and the
decidable kernel
\item The expansion of the propagator network's scope to include infinite
domains (through widening, coinduction, and abstract interpretation)
\end{enumerate}

Prior art: Agda's structural recursion checker (Level 2), Datalog's
stratified negation (stratified non-monotonicity), abstract
interpretation's widening operators (Level 3), Radul \& Sussman's
propagator model (Level 1 for finite domains). Prologos combines
these into a unified framework with explicit guarantee levels.
\section{References}
\label{sec:org59c6d89}

\begin{itemize}
\item Tarski, A. (1955). "A lattice-theoretical fixpoint theorem and its
applications." \emph{Pacific Journal of Mathematics}. (Level 1 foundation)
\item Cousot, P. \& Cousot, R. (1977). "Abstract Interpretation: A Unified
Lattice Model." (Levels 1, 3: widening/narrowing)
\item Abel, A. (2010). "MiniAgda: Integrating Sized and Dependent Types."
(Level 2: sized types as well-founded measures)
\item Radul, A. \& Sussman, G.J. (2009). "The Art of the Propagator."
(Level 1: monotone propagators on finite lattices)
\item Ceri, S., Gottlob, G., \& Tanca, L. (1989). "What You Always Wanted
to Know About Datalog." (Stratified fixpoints)
\item Turner, D.A. (2004). "Total Functional Programming." (Gödel
Completeness in functional languages)
\item Abel, A. \& Pientka, B. (2013). "Wellfounded Recursion with
Copatterns." (Level 4: coinductive productivity)
\end{itemize}

\begin{quote}
"A language where you can prove termination is more useful than one
where you can't --- not because non-termination is uncommon, but
because the \emph{ability to prove} it reveals structural understanding of
the computation."
\end{quote}
\end{document}
